% Options for packages loaded elsewhere
\PassOptionsToPackage{unicode}{hyperref}
\PassOptionsToPackage{hyphens}{url}
\documentclass[
  a4paper]{article}
\usepackage{xcolor}
\usepackage[margin=2.5cm]{geometry}
\usepackage{amsmath,amssymb}
\setcounter{secnumdepth}{-\maxdimen} % remove section numbering
\usepackage{iftex}
\ifPDFTeX
  \usepackage[T1]{fontenc}
  \usepackage[utf8]{inputenc}
  \usepackage{textcomp} % provide euro and other symbols
\else % if luatex or xetex
  \usepackage{unicode-math} % this also loads fontspec
  \defaultfontfeatures{Scale=MatchLowercase}
  \defaultfontfeatures[\rmfamily]{Ligatures=TeX,Scale=1}
\fi
\usepackage{lmodern}
\ifPDFTeX\else
  % xetex/luatex font selection
\fi
% Use upquote if available, for straight quotes in verbatim environments
\IfFileExists{upquote.sty}{\usepackage{upquote}}{}
\IfFileExists{microtype.sty}{% use microtype if available
  \usepackage[]{microtype}
  \UseMicrotypeSet[protrusion]{basicmath} % disable protrusion for tt fonts
}{}
\makeatletter
\@ifundefined{KOMAClassName}{% if non-KOMA class
  \IfFileExists{parskip.sty}{%
    \usepackage{parskip}
  }{% else
    \setlength{\parindent}{0pt}
    \setlength{\parskip}{6pt plus 2pt minus 1pt}}
}{% if KOMA class
  \KOMAoptions{parskip=half}}
\makeatother
\setlength{\emergencystretch}{3em} % prevent overfull lines
\providecommand{\tightlist}{%
  \setlength{\itemsep}{0pt}\setlength{\parskip}{0pt}}
\usepackage{bookmark}
\IfFileExists{xurl.sty}{\usepackage{xurl}}{} % add URL line breaks if available
\urlstyle{same}
\hypersetup{
  hidelinks,
  pdfcreator={LaTeX via pandoc}}

\author{}
\date{}

\begin{document}

\section{\texorpdfstring{Construction of \((R, Q)\) Mapping Adopting a
6-Dimensional
Hypercube}{Construction of (R, Q) Mapping Adopting a 6-Dimensional Hypercube}}\label{construction-of-r-q-mapping-adopting-a-6-dimensional-hypercube}

\textbf{Author}: Noriaki Kihara\\
\textbf{Affiliation}: WF System Co., Ltd.\\
\textbf{ORCID}: 0009-0004-6753-4020\\
\textbf{Version}: v1.0\\
\textbf{Date}: April 2026\\
\textbf{DOI}:
\href{https://doi.org/10.5281/zenodo.19692853}{10.5281/zenodo.19692853}

\begin{center}\rule{0.5\linewidth}{0.5pt}\end{center}

\subsection{Abstract}\label{abstract}

This paper defines an interface connecting the subjective spaces of the
three-layer structure model of central projection {[}M1{]} with the
combinatorial structure of a 6-dimensional hypercube {[}W8{]}.

This paper adopts two premises and, under those premises, defines
terminology and notation. It contains no theorem proofs. The purpose of
this paper is to provide a common definitional foundation that
subsequent papers may reference.

\textbf{Keywords}: 6-dimensional hypercube, central projection,
subjective space, metadata, \((R, Q)\) mapping

\begin{center}\rule{0.5\linewidth}{0.5pt}\end{center}

\subsection{§1 Adopted Premises}\label{adopted-premises}

This paper adopts the following two premises. These are not objects of
proof in this paper but are axiomatically adopted as the foundation for
the definitions constructed herein.

\subsubsection{Premise I (Adoption of the 6-Dimensional
Hypercube)}\label{premise-i-adoption-of-the-6-dimensional-hypercube}

We adopt the 6-dimensional hypercube defined in {[}W8{]} (a
combinatorial structure with axes \(x, y, z, t, R, Q\)). Here,
\(x, y, z, t, R\) are five axes taking real values, and
\(Q \in \{0, 1, 2, 3, 4, 5, 6, 7\}\) is an 8-valued discrete label axis
with 3-bit encoding \(Q = 4c_1 + 2c_2 + c_3\) (\(c_i \in \{0, 1\}\))
({[}W8{]} Definitions 1.1, 1.2).

\subsubsection{Premise II (Metadata-Holding
Capability)}\label{premise-ii-metadata-holding-capability}

A subjective space can hold arbitrary metadata that is numerically
representable.

\textbf{Remark}: Premise I is an external assumption of this paper and
does not claim that the 6-dimensional hypercube is derived from the
geometry of central projection (see {[}W8{]} §1). Premise II does not
restrict in advance the types of metadata that can be held.

\begin{center}\rule{0.5\linewidth}{0.5pt}\end{center}

\subsection{§2 Definitions}\label{definitions}

\subsubsection{Definition 2.1 (Schema and
Instance)}\label{definition-2.1-schema-and-instance}

The combinatorial structure of the 6-dimensional hypercube adopted in
Premise I is called the \textbf{schema} \(\mathcal{S}\). For each
subjective space \(S^n(R_1)\) (\(R_1 > 0\)), the concrete entity
obtained by assigning \(\mathcal{S}\) is called an \textbf{instance}.

\subsubsection{\texorpdfstring{Definition 2.2 (\((R, Q)\)
Mapping)}{Definition 2.2 ((R, Q) Mapping)}}\label{definition-2.2-r-q-mapping}

The map that identifies instances is called the \((R, Q)\) mapping and
is defined as follows:

\[
\varphi : \mathcal{I} \to \mathbb{R}_{> 0} \times \{0, 1, 2, 3, 4, 5, 6, 7\}
\]

Here, \(\mathcal{I}\) is the set of instances, and \(\varphi\) assigns
to each instance a pair consisting of a curvature radius
\(R_1 \in \mathbb{R}_{> 0}\) and a discrete label
\(Q \in \{0, 1, \ldots, 7\}\).

\textbf{Remark}: The \(R\) component is identified with the curvature
radius parameter \(R_1\) of {[}M1{]} §3.1. The \(Q\) component is the
value of the label axis of {[}W8{]}, held as metadata by virtue of
Premise II.

\subsubsection{Definition 2.3 (Uniqueness of Schema and Multiplicity of
Instances)}\label{definition-2.3-uniqueness-of-schema-and-multiplicity-of-instances}

The schema \(\mathcal{S}\) is unique. No restriction is imposed on the
cardinality of the instance set \(\mathcal{I}\). It is not excluded that
distinct instances may share the same \((R_1, Q)\) value.

\textbf{Remark}: The discrimination of multiple instances sharing the
same \((R_1, Q)\) is outside the definitional scope of this paper and is
deferred to subsequent papers.

\begin{center}\rule{0.5\linewidth}{0.5pt}\end{center}

\subsection{§3 Relationship with Prior
Works}\label{relationship-with-prior-works}

\subsubsection{§3.1 Relationship with the Three-Layer Structure Model
{[}M1{]}}\label{relationship-with-the-three-layer-structure-model-m1}

The instances in Definition 2.1 of this paper are a concretization of
the structure permitted by Condition 3 (non-constraint on multiplicity)
of {[}M1{]} §2. {[}M1{]} sets as a condition that infinitely many
subjective spaces may coexist at different curvature radii, and this
paper constructs its definitions within the scope of this condition.

\subsubsection{§3.2 Relationship with the 6-Dimensional Hypercube
{[}W8{]}}\label{relationship-with-the-6-dimensional-hypercube-w8}

This paper adopts the combinatorial structure of {[}W8{]} as the schema
\(\mathcal{S}\), but the 62-state enumeration results and the
interpretation examples of §9 of {[}W8{]} are not included in the
definitions of this paper. The internal structure of {[}W8{]} is
referenced independently as the content of the schema.

\begin{center}\rule{0.5\linewidth}{0.5pt}\end{center}

\subsection{§4 What This Paper Does Not
Claim}\label{what-this-paper-does-not-claim}

\begin{enumerate}
\def\labelenumi{\arabic{enumi}.}
\tightlist
\item
  That the 6-dimensional hypercube is isomorphic to the subjective
  spaces of central projection.
\item
  That the 6-dimensional hypercube is derived from the geometry of
  central projection.
\item
  That the definitions of this paper cannot be constructed with
  structures other than the 6-dimensional hypercube.
\item
  The physical existence of instances.
\item
  Physical interpretations of the \((R, Q)\) mapping (such as the micro
  black hole interpretation).
\item
  Methods for discriminating multiple instances sharing the same
  \((R_1, Q)\).
\end{enumerate}

\begin{center}\rule{0.5\linewidth}{0.5pt}\end{center}

\subsection{References}\label{references}

\begin{itemize}
\tightlist
\item
  {[}P1{]}: Noriaki Kihara, Geometric Formulation of Four-Dimensional
  Spacetime via Central Projection, Zenodo, 2026. DOI:
  \href{https://doi.org/10.5281/zenodo.19427780}{10.5281/zenodo.19427780}
\item
  {[}M1{]}: Noriaki Kihara, Three-Layer Structure Model of Central
  Projection: Nested Structure of Subjective Spaces via
  \(R\)--\(R_1\)--\(R_0\) and Middleware Geometry, Zenodo, 2026. DOI:
  \href{https://doi.org/10.5281/zenodo.19691713}{10.5281/zenodo.19691713}
\item
  {[}M2{]}: Noriaki Kihara, Regularity of Central Projection Between
  Hyperspherical Shells: A Formulation Without Tangent Hyperplanes and
  Degeneracy at \(R = 0\), Zenodo, 2026. DOI:
  \href{https://doi.org/10.5281/zenodo.19692192}{10.5281/zenodo.19692192}
\item
  {[}W8{]}: Noriaki Kihara, Set Structure and Combinatorial Properties
  of a 6-Dimensional Hypercube, Zenodo, 2026. DOI:
  \href{https://doi.org/10.5281/zenodo.19688521}{10.5281/zenodo.19688521}
\end{itemize}

\begin{center}\rule{0.5\linewidth}{0.5pt}\end{center}

\emph{v1.0}

\end{document}
